\section{Elicitation process}
The elicitation process consisted in three steps: (1) qualitative indicators, (2) value functions for quantitative indicators, and (3) \gls{pilebwt}.

\paragraph{Qualitative Indicators}
\label{sec:qualitative_indicators_uncertanty}
As presented in section \ref{sec:qi}, the author prepared an informed starting point for all qualitative indicators included in the analysis. The starting point determines the order of the alternatives and a linear value function, either increasing or decreasing, assuming all points are equally spaced. Hypothetical worst and hypothetical best points are included to give more flexibility during value function design. Experts are first requested to review the ranking and correct it if necessary. Then the linear value function can be adjusted to increase or decrease the difference between alternatives.


The indicators subject to this procedure are presented in table \ref{eq:qi_indicators} along with their descriptions. To improve the statistical robustness and given that this process is straightforward and not time-consuming, a total of nine experts were consulted (See Section \ref{sec:results} for details about the consulted experts).

\begin{table}[!ht]
    \centering
    \caption{Qualitative indicators elicited through ordering and value function adjustment.}
    \label{eq:qi_indicators}
    \begin{tabularx}{\textwidth}{lX}
        \toprule
        \textbf{Qualitative indicator} & \textbf{Description} \\
        \midrule
        Design maturity & How close to \gls{noak} installation \\
        Licensing status & How close / Ease of obtaining building licence (Country-dependant). \\
        Supplier availability & Ease of building up a supply chain within the country (Country-dependant). \\
        Design complexity & Proxy to determine likelihood of delays in the manufacturing and procurement of components. \\
        Construction complexity & Proxy to determine likelihood of delays in the construction and installation phases. (Country-dependant) \\
        \bottomrule
    \end{tabularx}
\end{table}

\paragraph{Value Functions}
\label{sec:value_functions}
The mid-value splitting technique, described in section \ref{sec:mavt_presentation}, is a well-established elicitation method for constructing value functions. However, its application is extremely time-consuming and cognitively demanding for experts. To address these challenges, this study carried out a \textit{first elicitation} adhering to the mid-value splitting technique to establish an initial set of value functions.

Subsequent elicitations employed a similar approach to that used for qualitative indicators (section \ref{sec:qualitative_indicators_uncertanty}), where the expert is first presented with an initial value function (the result of the \textit{first elicitation}) and is then asked to provide feedback and modify it as they see fit. This approach noticeably sped up the process, but is also likely to have introduced some anchoring bias. In regard of this study it is deemed not to be a major issue since the main goal is to evaluate the proposed methodology rather than obtaining definitive results.

\paragraph{Weights}
The last part of the elicitation is the definition of the weights applying the \gls{pilebwt} methodology presented in section \ref{sec:bwt}. This session was deemed intensive by the experts due to the intrinsic difficulty of the questions at hand. An observation that must not be interpreted as negative, since it properly reflects the difficulty of a trade-off problem. Nonetheless, this has to be considered by the analyst when planning elicitation sessions, by properly allocating time and breaks, to avoid tiredness and loss of focus that could compromise the quality of the answers. 

Unfortunately, in this case, it is not possible to circumvent the time-consuming nature of this task due to the method itself requiring the expert to reason about all trade-offs. To ensure consistency, since the \gls{bwt} method requires the knowledge of value functions, only experts who already participated in the value function elicitation sessions were involved in this step. 


